\chapter*{Hypotheses and Conventions Audit}
\phantomsection
\addcontentsline{toc}{chapter}{Hypotheses and Conventions Audit}

\section*{Branches and logarithms}
A branch of the logarithm is always understood on an explicitly specified
domain admitting a continuous argument.  Statements involving powers
\(z^\alpha=\exp(\alpha\operatorname{Log} z)\) inherit that branch choice.  Branch-cut
identities are not treated as global identities on \(\mathbb C^\times\).

\section*{Orientation and winding number}
Positively oriented simple closed contours are counterclockwise unless stated
otherwise.  The winding number is
\[
\operatorname{Ind}(\gamma,a)=\frac{1}{2\pi i}\int_\gamma\frac{dz}{z-a}.
\]
Residue-theorem signs follow this orientation convention.

\section*{Riemann-surface conventions}
Compactification, genus, divisor, and elliptic-function statements are to be
read with the hypotheses stated in the corresponding theorem.  Local
holomorphic coordinates are complex-oriented, and lattice quotients
\(\mathbb C/\Lambda\) use full rank-two lattices unless explicitly generalized.
